\documentclass[11pt]{article}
\usepackage{fontspec}
\usepackage[french]{babel}
\usepackage{amsmath}
\usepackage{amssymb}
\usepackage{amsthm}
\usepackage{mathrsfs}
\usepackage{tikz-cd}
\usepackage{enumitem}

\newcommand{\printedpage}[1]{\marginpar{\footnotesize p.~#1}}
\newcommand{\msp}[1]{\marginpar{\footnotesize [#1]}}

% footnote marks printed as superscript (n) in the text, plain n before the note
\makeatletter
\renewcommand\@makefnmark{\hbox{\textsuperscript{(\@thefnmark)}}}
\renewcommand\@makefntext[1]{\parindent 1em\noindent\hbox to 1.8em{\hss\textsuperscript{\@thefnmark}\,}#1}
\makeatother

\setlength{\parindent}{0pt}
\setlength{\parskip}{0.6em}

\begin{document}

\printedpage{44}
% head: A. Grothendieck : Dérivateurs, Ch. XVI, Édition des Universités de Montpellier II et Paris VII
% foot: Version 18 mars 2019

\msp{page 82}
de la démonstration de 5.5 b).)

Soit $x$ un objet (point) de $X$, $x_i$ son image dans $X_i$. Considérons (via ma construction standard des catégories de chemins) les $\Omega^1(X,x)$, $\Omega^1(X_i,x_i)$. On trouve
\[ \Omega^1(X,x) \simeq \varinjlim_{I} \Omega^1(X_i,x_i), \]
et on vérifie que c'est là encore un système inductif $\mathrm{L}_1$-$\pi_0$-adapté. En itérant, on trouve pour $n \in \mathbf{N}^*$
\[ \Omega^n(X,x) \simeq \varinjlim_{I} \Omega^n(X_i,x_i), \]
limite inductive $\mathrm{L}_1$-$\pi_0$-adaptée. Utilisant 5.6, on en conclut que pour tout $n$, le sous-ensemble (supposons $I$ un ensemble ordonné, que diable ça suffit !)
\[ I_n = \{ i \in I \mid \Omega^n(X_i,x_i) \ \text{ est 0-connexe} \} \]
est cofinal dans $I$, de plus il est stable par $\varinjlim$ filtrantes de cardinal $\leq \pi_0$. Donc % [sic] $\leq \pi_0$ -> $\leq \aleph_0$ (?), a cardinal bound is apparently intended
\[ I_\infty \overset{\text{déf}}{=} \bigcap I_n \]
est aussi cofinal dans $I$ (\textit{cf.} XVIII, 13.10.2, p. 302).

\msp{page 83}
Or $i \in I_\infty$ signifie que tous les $\Omega^n(X_i,x_i)$ sont 0-connexes, ce qui signifie aussi que $X_i$ est $\mathbf{W}_\infty$-asphérique (puisqu'il est déjà 0-connexe) (\textit{cf.} XII, cor. p. 190). Cela achève la démonstration du théorème 5.4.

\section*{6\quad La condition Loc (5, 7 D$_{W_\omega}$) implique $W \subset \mathbf{W}_\infty$}

Je vais redonner ici, pour mémoire, la démonstration du théorème 1 de XII :

\textbf{Théorème 6.1.} \textit{Soit $W$ un localiseur fondamental satisfaisant à} D$_{W_\omega}$ \textit{(où $W_\omega$ est le plus petit de tous les localiseurs fondamentaux). Alors on a $W = \mathrm{Fl}\,\mathrm{Cat}$,} ou \textit{$W \subset \mathbf{W}_\infty$ (suivant que $W$ ne satisfait pas à l'axiome de connexité} Loc (5), \textit{ou qu'il y satisfait.})

Je vais faire la démonstration en supposant que $W$ satisfait à Loc (5). Pour la démonstration du fait que dans le cas contraire, on a $W = \mathrm{Fl}(\mathrm{Cat})$, je renvoie à XII, th. 2, p. 20. (Voir aussi le résultat très voisin, plus général et plus faible, rappelé dans prop. 2.7, page 12.)

\textbf{Lemme 6.2.} \textit{Soient $W$, $W'$ deux localiseurs fondamentaux, $W$ satisfaisant à} Loc (7 D$_{W_\omega}$). \textit{Pour qu'on ait $W \subset W'$, il faut et il suffit qu'on ait} $\mathrm{Asph}_W \subset \mathrm{Asph}_{W'}$.

\msp{page 84}
Il suffit bien sûr de prouver “il suffit”, donc supposons que tout objet $W$-asphérique est $W'$-asphérique, et prouvons que $W \subset W'$. Soit donc $f : X \longrightarrow Y$ dans $W$, prouvons $f \in W'$. Pour ceci, j'utilise la factorisation canonique $f = pi$,
\[ X \xrightarrow{\ i \in W_\omega\ } \bar{X} \xrightarrow{\ p \in \mathrm{Parf}_{W_\omega}\ } Y. \]
\printedpage{45}
% head: A. Grothendieck : Dérivateurs, Ch. XVI, Édition des Universités de Montpellier II et Paris VII
% foot: Version 18 mars 2019
On a donc aussi $i \in W$, et comme $f = pi \in W$, on a $p \in W$, donc $p \in \mathrm{Parf}_{W_\omega} \cap W$. Or D$_{W_\omega}$ signifie que cet ensemble est contenu dans $W^{\mathrm{u}}$, donc on a $p \in W^{\mathrm{u}}$, et en particulier les fibres de $p$ sont $W$-asphériques. Elles sont donc $W'$-asphériques, et comme $p$ est $W_\omega$-parfait et a fortiori $W'$-parfait, il s'ensuit que $p \in W'$. On a aussi $i \in W_\omega \subset W'$, donc $i \in W'$, donc aussi $f = pi \in W'$, q.e.d.

Revenons alors à la démonstration de 6.1. En vertu du lemme, il suffit de prouver que si $X$ dans Cat est $W$-asphérique, il est $\mathbf{W}_\infty$-asphérique. Mais par hypothèse Loc (5) on sait déjà que $X$ est 0-connexe, et choisissant un point $x$ dans $X$, qu'il soit $\mathbf{W}_\infty$-asphérique s'exprime par la condition
\msp{page 85}
que les $\Omega^n(X,x)$ (pour $n \geq 1$) sont 0-connexes (\textit{cf.} XII, cor. p. 190 déjà cité). Pour ceci, à cause de Loc (5) il suffit de prouver qu'ils sont $W$-asphériques, et par induction il suffit de prouver que $\Omega^1(X,x)$ l'est. Or par définition de $\Omega^1(X,x)$, c'est là la fibre en $x$ d'une flèche dans Cat
\[ E = \underline{\mathrm{Ch}}(X,x,-) \xrightarrow{\ p \in \mathrm{Parf}_{W_\omega}\ } X, \]
où la source $E$ est $W_\omega$-asphérique, a fortiori $W$-asphérique. Comme $X$ l'est également, on a donc $p \in W$, donc $p \in \mathrm{Parf}_{W_\omega} \cap W \subset W^{\mathrm{u}}$, donc $p$ est à fibres $W$-asphériques, et en particulier $\Omega^1(X,x)$ est $W$-asphérique, q.e.d.

\textbf{Corollaire 6.3.} (Caractérisation de $\mathbf{W}_\infty$.) \textit{$\mathbf{W}_\infty$ est le plus grand localiseur fondamental satisfaisant à} Loc (5, 7 D$_{W_\omega}$) \textit{- donc aussi le plus grand localiseur fondamental satisfaisant à} Loc (5, 7).

Je pense que nous trouverons plus bas une caractérisation “interne” de $\mathbf{W}_\infty$, via un “axiome de descente” (très fort) Loc (10),
\msp{page 86}
comme l'unique localiseur fondamental satisfaisant Loc (5, 7 D$_{W_\omega}$, 10) (ou encore Loc (5, 7, 10)).

6.4. \textit{Commentaire sur les axiomes des $W$-fibrations du types} Loc (7 \ldots), \textit{et notamment sur} Loc (7 D$_{W_\omega}$). \textit{Relation conjecturale avec les} “types d'homotopie simples” \textit{de} \textsc{Whitehead}. % [sic] du types -> du type (or des types), grammatical agreement

a) Il me semble que \textit{tous} les localiseurs fondamentaux avec lesquels on a travaillé tacitement jusqu'à présent (les $\mathbf{W}_n$ ($n \in \mathbf{N}^*$), équivalence homologique et cohomologique, équivalence rationnelle \ldots) correspondent à des localiseurs $W \supset \mathbf{W}_\infty$. Par le théorème 6.1, il s'ensuit donc qu'\textit{aucun} d'eux, à part $\mathbf{W}_\infty$ lui-même, ne satisfait à la plus faible version D$_{W_\omega}$ de toutes de l'axiome Loc (7 \ldots), à la seule exception de Loc (7 F). Ainsi, les localiseurs $W$ qui y satisfont sont à l'exception de D$_{W_\omega}$ lui-même, d'un type tout nouveau - des localiseurs qui sont strictement contenus dans $\mathbf{W}_\infty$.

b) Ceci soulève, bien sûr, la question de l'existence de localiseurs fondamentaux contenus dans $\mathbf{W}_\infty$ et
\msp{page 87}
distincts de $\mathbf{W}_\infty$. Je n'ai pas même prouvé que $W_\omega \neq \mathbf{W}_\infty$, alors que je n'ai pas le moindre doute que tel est bien le cas. \texttt{[Cisinski a démontré que} $W_\omega = \mathbf{W}_\infty$\texttt{.]} A fortiori, j'ignore si le plus petit localiseur fondamental satisfaisant Loc (7 D$_{W_\omega}$) est différent de $\mathbf{W}_\infty$. Je serais tout particulièrement intéressé par le plus petit localiseur $\tilde{\mathbf{W}}_\infty \subset \mathbf{W}_\infty$ satisfaisant tous les axiomes de stabilité Loc (1 ter, 2, 3 bis, 4, 6, 7, 8) (\textbf{N.B.} il satisfait
\printedpage{46}
% head: A. Grothendieck : Dérivateurs, Ch. XVI, Édition des Universités de Montpellier II et Paris VII
% foot: Version 18 mars 2019
nécessairement Loc (5) car $\tilde{\mathbf{W}}_\infty \subset \mathbf{W}_\infty \subset \mathbf{W}_0$, de plus Loc (7) équivaut, en présence de Loc (8), à la forme très faible Loc (7 D$_{W_\omega}$) : $\mathrm{Parf}_{W_\omega} \cap W \subset W^{\mathrm{u}}$. Enfin, on montre que Loc (4) est conséquence des autres axiomes, \textit{cf.} plus bas.) Je conjecture à présent que $\tilde{\mathbf{W}}_\infty$ (qui donne lieu à un formalisme homotopique tout aussi complet que la sempiternelle “équivalence faible”) est bel et bien différent de $\mathbf{W}_\infty$ \texttt{[non par Cisinski]}. Et je suspecte
\msp{page 88}
que la catégorie $\mathrm{Hot}_{\tilde{\mathbf{W}}_\infty}$ mérite d'être appelée catégorie des \textit{“types d'homotopie simples”} (à la \textsc{Whitehead})\footnote[41]{Canulé (?), \textit{cf.} commentaire page 92.}. Un premier test décisif, c'est de vérifier que l'invariant cohomologique principal de \textsc{Whitehead} d'un type d'homotopie simple, savoir un élément dans un $K(\mathbf{Z}[\pi_1])$ (si je me rappelle bien\footnote[42]{non, c'est un invariant plus fin \ldots}) - ou quelque chose d'approchant, peut se définir déjà en termes des objets dans $\mathrm{Hot}_{\bullet\,\tilde{\mathbf{W}}_\infty}$ (le Hot pointé). Et deuxio, si on associe à une “maquette simpliciale” finie (décrivant combinatoirement une triangulation finie) l'ensemble ordonné de ses simplexes, vu comme un objet de Cat, et l'objet correspondant dans $\mathrm{Hot}_{\tilde{\mathbf{W}}_\infty}$, celui-ci ne change pas\footnote[43]{C'est évident même pour $W_\omega$, puisque le type d'homotopie \textit{stricte ne change pas !} - Non, il y a ici confusion entre homotopie topologique et homotopie combinatoire, \textit{cf.} p. 92 \ldots}, à isomorphisme canonique près, quand on effectue sur ce “polyèdre” envisagé une opération élé\msp{page 89}mentaire à la \textsc{Whitehead}. Cela suggère d'ailleurs d'essayer de définir de telles opérations élémentaires dans Cat tout entier, ou tout au moins dans la sous-catégorie des (petits) ensembles ordonnés, dans l'espoir de parvenir (peut-être) à une construction plus directe du localiseur pertinent pour le type d'homotopie simple.

c) Je termine par un commentaire sur le sens de l'axiome Loc (7 D$_{W_\omega}$), par opposition avec l'axiome plus faible Loc (7 F) qui assure seulement l'existence d'une “structure à fibrations” $(W, \mathrm{Fib}_W)$ sur Cat. Je signale aussi en passant que Loc (7 F), tout comme Loc (5) et comme Loc (9), et à la différence de tous les autres axiomes passés en revue jusqu'à présent, n'est \textit{pas} un axiome de stabilité : il affirme, non que certaines flèches (sous des hypothèses convenables) sont dans $W$, mais l'existence d'une factorisation. Aussi il n'est
\msp{page 90}
nullement clair que l'intersection de deux localiseurs fondamentaux satisfaisant à Loc (7 F) y satisfasse.

Loc (7 D$_{W_\omega}$) est exactement ce qu'il faut pour assurer que la factorisation canonique d'une flèche $f : X \longrightarrow Y$ de Cat en $f = pi$, avec $i \in W_\omega$, $p \in \mathrm{Parf}_{W_\omega}$, est aussi une factorisation du type stipulé par l'axiome de factorisation Loc (7 F) : on a (bien sûr) $i \in W$, mais aussi $p \in \mathrm{Fib}_W$ (car D$_{W_\omega}$ équivaut à sa forme forte $\mathrm{Parf}_{W_\omega} \subset \mathrm{Fib}_W$). Cela signifie aussi que la factorisation canonique de $f$ peut être utilisée pour “calculer” ses fibres homotopiques, non seulement en les points de $Y$, mais aussi par rapport à un morphisme de changement de base $Y' \longrightarrow Y$.

En l'absence de cet axiome,
\msp{page 91}
\printedpage{47}
% head: A. Grothendieck : Dérivateurs, Ch. XVI, Édition des Universités de Montpellier II et Paris VII
% foot: Version 18 mars 2019
admettant seulement Loc (7 F), on sait qu'il existe une théorie des fibres $W$-homotopiques, mais on n'a a priori aucun moyen pratique pour “calculer” lesdites fibres, faute de disposer d'un critère maniable suffisant pour assurer $p \in \mathrm{Fib}_W$. J'ai donc l'impression que l'axiome Loc (7 F) va être d'un maigre usage. (Il est vrai que l'exemple de $\mathbf{W}_0$ devrait me rendre prudent ! Il me faudrait regarder le cas des $\mathbf{W}_n$, s'ils satisfont Loc (7 F) \ldots). À mon sens, l'axiome Loc (7 D$_{W_\omega}$) est l'exact “minimum vital” pour être à l'aise dans le maniement des fibres $W$-homotopiques. Il est remarquable que ceci implique qu'on doive avoir $W \subset \mathbf{W}_\infty$, et que cela exclut du même coup \textit{toutes} les variantes considérées à ce jour de la notion de type d'homotopie !

\msp{page 92}
6.5. Finalement, la relation entrevue avec les “types d'homotopie simples” de \textsc{Whitehead} s'évanouit\footnote[44]{Mais peut-être quand même pas - la question reste ouverte.} : en effet, ces derniers sont strictement plus fin que les types d'homotopie (\textit{stricts}) ordinaires, et l'invariant de Whitehead n'est \textit{pas} un invariant du type d'homotopie ordinaire. Il est vrai que lorsqu'on considère le squelette simplicial d'un polyèdre triangulé comme objet de Cat, l'homotopie (continue) ordinaire entre deux tels polyèdres s'interprète, dans Cat, non pas (me semble-t-il) par l'existence d'un homotopisme entre les objets correspondants de Cat, mais par leur isomorphie dans $\mathrm{Hot}_{\mathbf{W}_\infty} \overset{\text{déf}}{=} \mathrm{Cat}(\mathbf{W}_\infty)^{-1}$. Mais le procédé d'approximations simpliciales de \textsc{Brouwer} montre que si les polyèdres $X$ et $Y$ sont homotopes, alors il existe des subdivisions simpliciales $X'$ et $Y'$ de $X$ et de $Y$ qui sont “combinatoirement homotopes” (dans un sens convenable). Donc je présume que les maquettes simpliciales de $X'$ et de $Y'$ sont alors homotopes dans Cat. Il s'ensuivrait que % [sic] plus fin -> plus fins, agreement with "ces derniers"
\msp{page 93}
$X$ et $Y$, j'entends leurs maquettes, sont bien homotopes dans Cat, s'il est vrai qu'une subdivision simpliciale (il suffirait la subdivision barycentrique, en itérant \ldots) $X'$ de $X$ a une maquette dans Cat qui lui est homotope. La question soulevée ici me paraît intéressante, il faut visiblement en trouver la réponse. Elle se pose déjà pour tout ensemble ordonnée $I$, qu'on va supposer fini pour simplifier, en prenant l'ensemble (ordonné par inclusion) $I' = \mathrm{sd}(I)$ des parties totalement ordonnées de $I$. On sait que les réalisations topologiques de $I$ et de $\mathrm{sd}(I)$ sont canoniquement isomorphes, mais $I$ et $\mathrm{sd}(I)$ sont-ils homotopes dans Cat ? On a bien une application croissante canonique % [sic] ensemble ordonnée -> ensemble ordonné, agreement
\[ f : \mathrm{sd}(I) \longrightarrow I, \]
qui associe à toute partie totalement ordonnée de $I$ son plus grand élément. Mais existe-t-il une application croissante en sens inverse, qui soit un homotopisme inverse du précédent à homotopie près ? Je viens de vérifier dans le cas

\msp{page 94} manque

\msp{page 95}
du bord du 2-simplexe

\begin{center}
\begin{tikzpicture}[baseline=(b.base)]
\node (b) at (0.95,0.3) {};
\draw (0,0) -- (1.9,0) -- (0.95,0.95) -- cycle;
\fill (0,0) circle (2.2pt);
\fill (1.9,0) circle (2.2pt);
\fill (0.95,0.95) circle (2.2pt);
\end{tikzpicture}
\qquad\qquad\qquad ou \qquad\qquad\qquad
\begin{tikzpicture}[baseline=(c.base)]
\node (c) at (0,-0.1) {};
\draw (0,0) circle (0.48);
\fill (90:0.48) circle (2.2pt);
\fill (210:0.48) circle (2.2pt);
\fill (330:0.48) circle (2.2pt);
\end{tikzpicture}
\qquad\quad ,
\end{center}

subdivision barycentrique
\printedpage{48}
% head: A. Grothendieck : Dérivateurs, Ch. XVI, Édition des Universités de Montpellier II et Paris VII
% foot: Version 18 mars 2019

\begin{center}
\begin{tikzpicture}[baseline=(b.base)]
\node (b) at (0.95,0.3) {};
\draw (0,0) -- (1.9,0) -- (0.95,0.95) -- cycle;
\fill (0,0) circle (2.2pt);
\fill (0.95,0) circle (2.2pt);
\fill (1.9,0) circle (2.2pt);
\fill (0.475,0.475) circle (2.2pt);
\fill (1.425,0.475) circle (2.2pt);
\fill (0.95,0.95) circle (2.2pt);
\end{tikzpicture}
\qquad\qquad\qquad ou \qquad\qquad\qquad
\begin{tikzpicture}[baseline=(c.base)]
\node (c) at (0,-0.1) {};
\draw (0,0) circle (0.48);
\fill (30:0.48) circle (2.2pt);
\fill (90:0.48) circle (2.2pt);
\fill (150:0.48) circle (2.2pt);
\fill (210:0.48) circle (2.2pt);
\fill (270:0.48) circle (2.2pt);
\fill (330:0.48) circle (2.2pt);
\end{tikzpicture}
\end{center}

(hexagone), que pour toute application croissante $g : I \longrightarrow \mathrm{sd}(I)$, sa réalisation topologique est homotope à une application constante, donc $g$ ne peut être un homotopisme dans Cat (car cela impliquerait que sa réalisation topologique est un homotopisme).

Ceci montre donc que la notion d'homotopisme dans Cat est \textit{strictement plus fine} que celle pour les réalisations topologiques, ce qui redonne à nouveau substance à l'espoir que l'invariant de Whitehead pourrait se définir comme invariant d'homotopie sur Cat, ou du moins sur une sous-catégorie pleine convenable, telle celle des ensembles ordonnés (éventuellement restreints à être finis), ou des maquettes simpliciales. Et même, que l'ensemble des Wh-types d'homotopie (types d'homotopie simples) peut se
\msp{page 96}
plonger canoniquement dans un $\mathrm{Hot}_W$ convenable, où $W$ serait un localiseur fondamental strictement plus fin que $\mathbf{W}_\infty$, p.ex. (qui sait ?) $\tilde{\mathbf{W}}_\infty$. Le premier travail qu'on aurait envie de faire, c'est de construire un tel $W$ \textit{ad hoc}, pour que l'invariant de Whitehead puisse s'écrire comme un foncteur sur $\mathrm{Hot}_{\bullet W}$. Pour commencer, il faudrait déjà s'assurer que cet invariant peut être vu comme un invariant du type d'homotopie, \textit{i.e.} qu'il s'interprète comme un foncteur sur $\mathrm{Hot}_{\bullet}\mathrm{Cat} \simeq \mathrm{Cat}(\,\underbrace{\text{homotop}}_{\substack{\text{ensemble des}\\ \text{homotopismes}}}\,)^{-1}$. Ce travail suggérera peut-être aussi la définition des “invariants de Whitehead supérieurs” (impliquant les $\pi_i$, non seulement le $\pi_1$), lesquels pourraient peut-être servir (avec celui de dimension 1) à décrire le localiseur hypothétique $W$ (“de Whitehead”).

\msp{page 97}

\section*{7\quad Commentaires sur les axiomes}

7.1. La principale raison d'être d'un localiseur fondamental $W$, c'est de définir une théorie d'homotopie’, et plus précisément un prédérivateur $\mathrm{HOT}_W$, dont on espère que ce sera même un dérivateur, ou qu'il satisfasse du moins à certains des axiomes des dérivateurs. Les axiomes Loc (1) à Loc (9) sont en majeure partie orientés vers cet objectif, à cela près que Loc (1, 2, 3) constitue un “minimum vital” absolu, pour pouvoir commencer à travailler avec $W$. Les principales questions qui se posent me paraissent les suivantes : % [sic] d'homotopie’ -> d'homotopie (or “théorie d'homotopie”), closing quote without opening quote
\begin{enumerate}[label={}, leftmargin=2.5em]
\item[1)] Quand Cat munie de $(W, \mathrm{Cof}_W)$, resp. de $(W, \mathrm{Fib}_W)$, est-elle une catégorie à cofibrations resp. à fibrations, de façon que les résultats de K. S. \textsc{Brown} soient applicables ?
\item[2)] Sous quelles conditions, sur $W$ et sur $f$, $f_!$ et $f_*$ existent-ils ? Et quand a-t-on, de plus, exactitude des carrés cocartésiens resp. cartésiens ?

\msp{page 98}
\item[3)] Plus généralement, en termes des axiomes imposés à $W$, quelles sont les axiomes Der (\ldots) que sont vérifiés pour le prédérivateur $\mathrm{HOT}_W$ ? (Ce qui pourrait aussi dépendre du choix du domaine Diag $\subset$ Cat. Le prédérivateur est défini a priori sur Cat tout entier). En particulier, % [sic] quelles sont les axiomes -> quels sont les axiomes, agreement; [sic] que sont vérifiés -> qui sont vérifiés, relative pronoun
\end{enumerate}

\end{document}
